\section*{Supplementary Note 5 | Chemical-overlap audit and recovery sensitivity}

We quantified chemical overlap between benchmark reference actives, the fine-tuning training partition and the ligand-only pretraining library to establish how much of the retrospective ECFP4-Tanimoto recovery could be explained by chemical familiarity rather than by pocket-conditioned generation. Per-molecule counts and percentages are in Supplementary Table~5, and recovery values under each exclusion are in Extended Data Tables~6 and~7; this note summarises what they show.

Identity between NA2smol outputs and fine-tuning training ligands is rare, so the fine-tuning split does not explain the recovery signal. Higher similarity to the pretraining library is expected, as pretraining exposes the model to a broad chemical space, and targets with reference actives in that space are the ones the model is designed to handle. Excluding targets with actives resembling the fine-tuning partition leaves NA2smol ahead on both benchmarks, confirming the fine-tuning split does not drive the ranking. Excluding targets with actives resembling the pretraining library changes the order: NA2smol falls to 0.053 on RNA (ED2Mol: 0.085) and to 0.238 on DUD-E (DiffSBDD: 0.455). DiffSBDD's DUD-E recovery remains nearly unaffected (0.447 to 0.455). These shifts follow from removing targets that the pretraining stage was designed to cover.

Excluding reference actives by structure identifier rather than by record yields the same results: the eligible-target set shrinks, but the recovery comparison between methods remains unchanged. These sensitivity analyses support the reported benchmark.
